\section{Forward heat, character dynamics, and a rational flag limit}

\begin{theorem}[NN-A015: a universal normalized character limit]
Retain the exact heat normalization of NN-A007 and set, for $t\ge0$,
\[
 p_m(t)=\sum_{j\ge1}x_j(t)^{-2m},\qquad
 \beta_j(t)=\frac{x_j(t)^{-2}}{p_1(t)},\qquad
 E_t(w)=F_t(w/p_1(t)).
\]
Under A these are positive zero weights, indexed in decreasing order
with multiplicity, and $\sum_j\beta_j(t)=1$. They obey the exact
moment equations
\[
 p_m'=4m\sum_{j=1}^{m}p_jp_{m+1-j}-m(4m+2)p_{m+1}
 \qquad(m\ge1).                                      \tag{A015.1}
\]
The moments are analytic even at a possible multiple zero at $t=0$;
no separate differentiable labeling of colliding roots is needed.

As $t\to+\infty$,
\[
 p_1(t)\sim\frac{(\log t)^2}{32},\qquad
 E_t(w)\longrightarrow E_\infty(w)=\cosh\sqrt{2w}
                                                        \tag{A015.2}
\]
locally uniformly on $\mathbb C$. The square-root expression is its
entire even-power series. In addition,
\[
 \beta_j(t)\longrightarrow
 \beta_j^*=\frac{8}{\pi^2(2j-1)^2},\qquad
 \sum_j|\beta_j(t)-\beta_j^*|\longrightarrow0.           \tag{A015.3}
\]
Consequently the characters constructed in NN-A008 have the limits
\[
 \chi_{S,t}(\sigma)\longrightarrow
 \prod_{k\ge2}\left[
 (-1)^{k-1}\frac{2^k(4^k-1)\zeta(2k)}{\pi^{2k}}
 \right]^{m_k(\sigma)},                                \tag{A015.4}
\]
and, for $U\in U(d)$ with eigenvalues $\lambda_1,\ldots,\lambda_d$,
\[
 \chi_{U,t}(U)\longrightarrow
 \prod_{a=1}^d\cosh\sqrt{2(\lambda_a-1)}.                \tag{A015.5}
\]
The limit in (A015.2) and the asymptotic for $p_1=c_1$ follow from
the kernel without A. A identifies the whole path $t\ge0$ with
positive zero weights and the indicated characters.
\end{theorem}

\begin{proof}
For every fixed real $t$, Fourier decay gives $H_t(x)\to0$ as
$x\to\pm\infty$ on the real line. If its zero product were finite,
the order-less-than-one product for $F_t$ would make $H_t$ a nonzero
polynomial, contradicting this decay. Under A its real zero product
therefore has infinitely many positive parameters at each $t\ge0$.
The heat PDE, in the coordinate $w$, and normalization at $w=0$ give
\[
 \partial_tF_t=4wF_t''+2F_t'-2p_1F_t.
\]
Near $w=0$, put $q=F_t'/F_t$ and use the convergent product to write
\[
 \log F_t(w)=\sum_{m\ge1}\frac{(-1)^{m-1}}m p_m(t)w^m,
 \qquad q(w)=\sum_{m\ge0}(-1)^m p_{m+1}(t)w^m.
\]
Then $\partial_t\log F_t=4w(q'+q^2)+2q-2p_1$.
Comparison of each coefficient proves (A015.1). Every $p_m$ is a
fixed polynomial in $c_1,\ldots,c_m$ by this logarithmic expansion,
so it is analytic in $t$ even if zero branches meet. In particular,
\[
 p_1'=4p_1^2-6p_2>0.                                  \tag{A015.6}
\]
Strictness also follows directly from the nonzero variance below.

We give a direct concentration proof of (A015.2), including the
tail control needed for local uniform convergence. The kernel is
\[
 \Phi(u)=\sum_{n\ge1}
 (2\pi^2n^4e^{9u}-3\pi n^2e^{5u})e^{-\pi n^2e^{4u}}
 \qquad(u\ge0).
\]
There are constants $0<b<B<\infty$ such that, on this whole half-line,
\[
 b e^{9u-\pi e^{4u}}\le\Phi(u)
       \le B e^{9u-\pi e^{4u}}.                        \tag{A015.7}
\]
The first summand proves the lower bound, since
$2\pi^2-3\pi e^{-4u}\ge2\pi^2-3\pi>0$.
For the upper bound discard the negative term and bound the remaining
sum by $2\pi^2\sum_{n\ge1}n^4e^{-\pi(n^2-1)}$ after factoring
out $e^{9u-\pi e^{4u}}$.

Let $U_t\ge0$ have density $H_t(0)^{-1}e^{tu^2}\Phi(u)$, and put
$a_t=(\log t)/4$ for large $t$. Integrating on
$[a_t,a_t+1/t]$ in (A015.7) gives
\[
 H_t(0)\ge\exp(t a_t^2-Ct)                            \tag{A015.8}
\]
for an absolute $C$ and all sufficiently large $t$.
Fix $0<\epsilon<1$, an integer $p\ge0$, and $A_0\ge0$.
On $0\le u\le(1-\epsilon)a_t$, the extra weight
$(u/a_t)^p e^{A_0u/a_t}$ is bounded. The integral over this interval,
divided by (A015.8), is at most a constant times
\[
 a_t\exp\bigl(-[2\epsilon-\epsilon^2]t a_t^2
                         +Ct+9a_t+A_0\bigr)\longrightarrow0.
\]
On $u\ge b_t=(1+\epsilon)a_t$, for all sufficiently large $t$,
\[
 tu^2+9u+A_0u/a_t+p\log(u/a_t)\le\tfrac\pi2e^{4u}.
\]
Indeed $u^2e^{-4u}$ and $ue^{-4u}$ decrease for large $u$, and
$t b_t^2/e^{4b_t}=O((\log t)^2t^{-\epsilon})\to0$;
the remaining linear terms are smaller. The weighted tail integral
is therefore at most
\[
 B\int_{b_t}^\infty e^{-(\pi/2)e^{4u}}du
 \le \frac{B}{2\pi e^{4b_t}}
                 e^{-(\pi/2)e^{4b_t}},
\]
which still tends to zero after division by (A015.8).
These estimates prove $U_t/a_t\to1$ in every fixed $L^p$, and
uniform integrability of all fixed exponential weights in $U_t/a_t$.
Since $p_1=c_1=\mathbb E U_t^2/2$, the claimed asymptotic for $p_1$
follows. The variance identity
$p_1'=\operatorname{Var}(U_t^2)/2>0$ justifies (A015.6).

Set $V_t=U_t/\sqrt{p_1(t)}$. We have $V_t\to\sqrt2$ in every
fixed $L^p$ and the same exponential integrability. Moreover
\[
 E_t(w)=\mathbb E\cosh(\sqrt w V_t).
\]
For $|w|\le R$ the difference from $\cosh\sqrt{2w}$ is bounded by
\[
 \sqrt R\,\mathbb E\left[
 |V_t-\sqrt2|e^{\sqrt R(V_t+\sqrt2)}\right],
\]
which tends to zero by Cauchy--Schwarz, the $L^2$ convergence and the
exponential-moment bound. This proves local uniform convergence.

The positive zeros of
$H_t(y/\sqrt{p_1(t)})/H_t(0)=E_t(-y^2)$ converge, in their order,
to the simple positive zeros of $\cos(\sqrt2y)$, by Rouche's theorem
and uniform nonvanishing on the intervening compact real intervals.
Under A all zeros being counted are real. Thus
\[
 \sqrt{p_1(t)}x_j(t)\to\frac{\pi(2j-1)}{2\sqrt2},
\]
which gives pointwise convergence in (A015.3). Both the finite-time
and limiting sequences sum to one; the latter uses
$\sum_{j\ge1}(2j-1)^{-2}=\pi^2/8$. Pointwise convergence of
nonnegative unit-mass sequences then gives $\ell^1$ convergence:
first truncate at $J$, pass to the limit in that finite set, and
bound the two omitted tails by their remaining masses.

For every fixed $k$, $|a^k-b^k|\le k|a-b|$ on $[0,1]$, hence
the power sums converge. Summing the odd reciprocal powers gives
\[
 \sum_j(\beta_j^*)^k
 =\frac{8^k}{\pi^{2k}}(1-2^{-2k})\zeta(2k)
 =\frac{2^k(4^k-1)\zeta(2k)}{\pi^{2k}}.
\]
The finite-cycle product gives (A015.4). Finally
$\chi_{U,t}(U)=\prod_{a=1}^d E_t(\lambda_a-1)$, so (A015.5)
is the already proved local uniform limit. These are actual
positive-definite characters, either by the constructions in NN-A008
with $\beta^*$ or by the pointwise limit of their finite positive
Gram matrices.
\end{proof}

\begin{corollary}[Rational normalized flag coordinates and the classical limit]
Fix $n\ge1$ and form the $(n+1)$ by $(n+1)$ upper Toeplitz matrix
from the coefficients of $E_t(w)$. Let $\widetilde\Delta_i(t)$
be its top-right $i$ by $i$ minor, with boundary values
$\widetilde\Delta_0=\widetilde\Delta_{n+1}=1$, and let
\[
 \widetilde q_i(t)=
 \frac{\widetilde\Delta_{i-1}(t)\widetilde\Delta_{i+1}(t)}
 {\widetilde\Delta_i(t)^2}.
\]
Every limit $\widetilde q_i^*$ is a strictly positive rational number,
computed from the Toeplitz coefficients $2^k/(2k)!$. In particular,
\[
 \widetilde q_1^*=n(2n-1)(4n-3).
\]
For the unnormalized coefficients of $F_t$, the corresponding
Rietsch quantum flag coordinates satisfy
\[
 q_i(t)=p_1(t)^{-2}\widetilde q_i(t)
 \sim\frac{1024\widetilde q_i^*}{(\log t)^4}.
\]
Thus this specific positive path approaches the locus $q_1=\cdots=q_n=0$
of the quantum parameters, whose specialization is the classical
cohomology ring in the usual quantum-cohomology presentation.
\end{corollary}
\begin{proof}
Coefficient convergence gives minor convergence. The limiting product
has infinitely many strictly positive parameters $\beta_j^*$, so
NN-A008 makes every allowed limiting minor strictly positive. Its
coefficients are rational, proving the first assertion. For $n\ge2$,
put $a_k=2^k/(2k)!$. Then
\[
 \widetilde\Delta_1^*=a_n,\qquad
 \widetilde\Delta_2^*=a_{n-1}^2-a_{n-2}a_n,
\]
and substitution gives
\[
 \widetilde q_1^*=
 n^2(2n-1)^2-n(n-1)(2n-1)(2n-3)
 =n(2n-1)(4n-3).
\]
For $n=1$, both boundary minors are one and $a_1=1$, giving the same
formula. Each term of $\Delta_i$ has coefficient degree
$i(n+1-i)$ under the scaling $c_k\mapsto p_1^{-k}c_k$.
Taking the second difference of this degree gives exactly
$\widetilde q_i=p_1^2q_i$. The asymptotic follows from (A015.2).
The coordinate identification is the imported Rietsch Theorem~2.1,
equations~(2.1)--(2.3), and Section~2.1, already checked in NN-A008;
the $q=0$ specialization is that theorem's standard quantum flag
presentation, not a new claim of an arithmetic motive.
\end{proof}

The limiting symmetric-group cycle values are rational because
$\zeta(2k)/\pi^{2k}$ is rational. This supplies actual rational finite
flag points, while the limiting Thoma parameters themselves contain
$\pi^{-2}$. Neither fact gives an integral positive basis for a GCT
multiplicity problem. The moment evolution and the universal limit
belong to this specified analytic heat family; they are not a theorem
about arbitrary $L$-function deformations.
